% Propietario conservado: /Users/ruben/Documents/New project/output/EXTREMA_MEDIA_RAZON_RECIPROCIDAD_ES_EN_20260918/ARTICULO/ES/source/sections/01_hilbert_nonadico.tex
% Procedencia: integral 2249, cap. 23, pp. 601--608; propietario base_c21_sin_encabezado.tex.
% Recuperación y desarrollo de pruebas. Requiere el núcleo común APP--TRIT--TPK.
\section{Réalisation hilbertienne de la phase nonadique et conservation de l’unité}
\label{x_reciprocidad:iv:hilbert}

La construction de l'holographie modulaire triadique fournit d'abord les états, leurs transitions et les applications de restriction. La réalisation hilbertienne ajoute à leurs coordonnées discrètes une structure linéaire et un produit scalaire. Cet ordre détermine ce que représente chaque opérateur : les étiquettes de la base proviennent de la construction nonadique ; les opérations linéaires agissent ensuite sur ces étiquettes. La mémoire généalogique demeure dans l'état du TPK et accompagne les projections utilisées pour le représenter.

\subsection{Fibre de phase, base étiquetée et produit scalaire}

Pour une profondeur entière \(d\geq1\), soit
\begin{equation}
 \mathcal C_{9,d}=(\ZZ/9\ZZ)^d,
 \qquad \mathcal H_d=\ell^2(\mathcal C_{9,d}),
 \qquad \mathcal A_d=B(\mathcal H_d).
 \label{x_reciprocidad:iv:eq:hilbert-finito}
\end{equation}
Le produit scalaire est
\(\langle f,g\rangle=\sum_x\overline{f(x)}g(x)\). Les fonctions
\(\delta_x\) forment une base orthonormale étiquetée par les coordonnées
de phase. Ainsi, \(\dim\mathcal H_d=9^d\) et
\(\mathcal A_d\cong M_{9^d}(\CC)\). L’algèbre matricielle représente
cette fibre finie ; les routes, les quotients et la mémoire qu’une
projection de phase omet font toujours partie de l’état de départ.

La factorisation des coordonnées induit l’unitaire
\begin{equation}
 \mathcal H_{d+1}\longrightarrow\mathcal H_d\otimes\CC^9,
 \qquad \delta_{(x,r)}\longmapsto\delta_x\otimes\delta_r.
 \label{x_reciprocidad:iv:eq:factorizacion-hilbert}
\end{equation}
Nous définissons la trace normalisée
\(\tau_d(a)=9^{-d}\operatorname{Tr}(a)\). Le dénombrement fini et la
normalisation se distinguent : \(\operatorname{Tr}(I)=9^d\), tandis
que \(\tau_d(I)=1\). De même, le groupe \((\ZZ/9\ZZ)^3\) et l’espace
vectoriel \(\FF_3^6\) ont tous deux \(729\) éléments, mais
conservent des lois algébriques différentes. La réalisation d’une fibre de
phase ne remplace pas l’encodage ternaire des émissions APP.

\subsection{Translation et phase : représentation finie de Weyl}

Dans la réalisation complexe postérieure aux constructions de \(\pi\) et de la structure \(\mathrm i^2=-1\), nous fixons \(\omega=\exp(2\pi\mathrm i/9)\). Sur \(\mathcal H_1\), soient
\begin{equation}
 (Xf)(r)=f(r-1),\qquad (Zf)(r)=\omega^r f(r).
 \label{x_reciprocidad:iv:eq:weyl-operadores}
\end{equation}
La première opération déplace la coordonnée de phase ; la seconde
l’évalue au moyen d’un caractère. Toutes deux sont unitaires et vérifient
\begin{equation}
 X^9=Z^9=I,\qquad ZX=\omega XZ.
 \label{x_reciprocidad:iv:eq:weyl-relacion}
\end{equation}
En effet, \((ZXf)(r)=\omega^r f(r-1)\), tandis que \((XZf)(r)=\omega^{r-1}f(r-1)\). Le caractère apparaît ici comme une réalisation de la phase déjà construite et ne sélectionne pas les registres numériques du TPK.

% RESULTADO_RECUPERADO: 14_numero_heisenberg_y_d108.tex:741--788.
% Prueba algebraica desarrollada con la convención X^a Z^b ya fijada aquí.
\subsubsection{Aire orientée et extension centrale des translations}
\label{x_reciprocidad:iv:weyl-area-extension}

Le support des translations est \(X_9=(\ZZ/9\ZZ)^2\).
La courbure mixte d’APP construite dans la
section~\ref{x_reciprocidad:iv:extension-curvatura-mixta} attribue \(+1\) à la
plaquette orientée de la paire somme--produit et \(-1\) à la
paire d’ordre contraire. Son extension bilinéaire à deux déplacements
\(u=(a,b)\) et \(v=(c,d)\) est la forme alternée
\begin{equation}
 \sigma_{\mathrm{ar}}(u,v)=ad-bc\quad\text{dans }\ZZ/9\ZZ.
 \label{x_reciprocidad:iv:eq:weyl-area-orientada}
\end{equation}
En effet, l’alternance annule l’évaluation sur deux déplacements
du même axe. Les valeurs sur la base orientée sont
\(\sigma_{\mathrm{ar}}((1,0),(0,1))=1\) et
\(\sigma_{\mathrm{ar}}((0,1),(1,0))=-1\) ; distribuer les deux
arguments donne \eqref{x_reciprocidad:iv:eq:weyl-area-orientada}. La même formule est
la somme des plaquettes de toute chaîne cellulaire orientée qui réalise
ce parallélogramme : les arêtes intérieures s’annulent. Modifier un
représentant entier par un multiple de neuf conserve sa publication
dans \(\ZZ/9\ZZ\).

L'ordre \(X^aZ^b\) des opérateurs déjà définis fixe une section de l'extension centrale. Notons son facteur de composition
\begin{equation}
 c_{\mathrm W}(u,v)=bc\quad\text{dans }\ZZ/9\ZZ.
 \label{x_reciprocidad:iv:eq:weyl-cociclo-polarizado}
\end{equation}
Ce cocycle conserve l'ordre des deux opérations, tandis que \(\sigma_{\mathrm{ar}}\) conserve l'orientation du parallélogramme.

\begin{proposition}[Extension centrale et signe du commutateur]
\label{x_reciprocidad:iv:prop:heisenberg-orientado}
L’ensemble \(\mathfrak H_9=(\ZZ/9\ZZ)^3\), muni du produit
\begin{equation}
 (a,b,t)(c,d,s)=(a+c,b+d,t+s+bc),
 \label{x_reciprocidad:iv:eq:heisenberg-ley}
\end{equation}
est un groupe d'ordre \(729\). L'application
\[
 \varrho(a,b,t)=\omega^tX^aZ^b
\]
est une représentation unitaire fidèle sur \(\mathcal H_1\).
Avec la convention \([g,h]=ghg^{-1}h^{-1}\), on a
\begin{equation}
 [\varrho(a,b,0),\varrho(c,d,0)]
 =\omega^{bc-ad}I
 =\omega^{-\sigma_{\mathrm{ar}}((a,b),(c,d))}I.
 \label{x_reciprocidad:iv:eq:weyl-conmutador-orientado}
\end{equation}
Le centre et le groupe des translations forment la suite exacte
\[
 0\longrightarrow\ZZ/9\ZZ
 \xrightarrow{\ t\mapsto(0,0,t)\ }\mathfrak H_9
 \xrightarrow{\ (a,b,t)\mapsto(a,b)\ }X_9
 \longrightarrow0.
\]
\end{proposition}
\begin{proof}
Pour \(w=(e,f)\), l’identité de cocycle s’écrit
\[
 c_{\mathrm W}(u,v)+c_{\mathrm W}(u+v,w)
 =bc+(b+d)e
 =de+b(c+e)
 =c_{\mathrm W}(v,w)+c_{\mathrm W}(u,v+w).
\]
Cela prouve l’associativité de \eqref{x_reciprocidad:iv:eq:heisenberg-ley} ; l’identité
est \((0,0,0)\) et l’inverse de \((a,b,t)\) est
\((-a,-b,-t+ab)\). Les trois résidus sont libres, d’où l’ordre
\(9^3=729\). De \(ZX=\omega XZ\), il résulte
\[
 (X^aZ^b)(X^cZ^d)=\omega^{bc}X^{a+c}Z^{b+d},
\]
qui démontre la propriété de représentation. Si \(\omega^tX^aZ^b=I\), son action sur \(\delta_r\) impose d'abord \(a=0\). L'évaluation en \(r=0\) donne \(t=0\), et l'évaluation en \(r=1\) donne \(b=0\), puisque \(\omega\) est d'ordre neuf. La représentation est donc fidèle. La comparaison des produits dans les deux ordres donne le facteur \(bc-ad\). Un élément central doit satisfaire \(bc-ad=0\) pour tout \((c,d)\). En prenant successivement \((c,d)=(1,0),(0,1)\), on obtient \(b=a=0\). Cela identifie le centre et prouve l'exactitude.
\end{proof}

La phase de Wilson pour l’orientation
\eqref{x_reciprocidad:iv:eq:weyl-area-orientada} est
\(\omega^{\sigma_{\mathrm{ar}}(u,v)}\) ; le commutateur de la
section ordonnée \(X^aZ^b\) publie son inverse. En particulier,
\(c_{\mathrm W}(u,v)-c_{\mathrm W}(v,u)
=-\sigma_{\mathrm{ar}}(u,v)\).
Cette identité relie deux conventions déterminées, sans identifier
le cocycle polarisé à la forme alternée. L’extension appartient
au groupe des translations du support APP. Sa représentation agit sur
la fibre de phase ; le registre de la trajectoire et les retenues conservent
leur résidence dans l’état TPK qui a produit cette fibre. L’ordre \(729\)
du groupe et le cardinal \(729\) de l’espace ambiant ternaire conservent les
types distincts établis dans \eqref{x_reciprocidad:iv:eq:hilbert-finito}.
Le centre \(\ZZ/9\ZZ\) de cette extension enregistre la phase
du commutateur. La mémoire \(C_{12}\) de la suite
\(0\to C_{12}\to C_{108}\to C_9\to0\), construite dans la
proposition~\ref{x_reciprocidad:prop:k-extension}, appartient au calendrier
dodécaphasique. Chaque extension conserve sa loi de composition et
sa projection propres.

\begin{proposition}[Algèbre engendrée par la phase et la translation]
\label{x_reciprocidad:iv:prop:weyl-completo}
Les \(81\) opérateurs \(X^aZ^b\), avec \(0\leq a,b<9\), forment
une base de \(\mathcal A_1=M_9(\CC)\). En particulier, la
représentation de \eqref{x_reciprocidad:iv:eq:weyl-operadores} est irréductible.
\end{proposition}
\begin{proof}
Si \(a\neq0\pmod9\), la matrice \(X^aZ^b\) n’a pas de coefficients
diagonaux. Si \(a=0\), sa trace est
\(\sum_{r=0}^8\omega^{br}\), égale à \(9\) pour \(b=0\) et à
\(0\) dans les autres cas. Par \eqref{x_reciprocidad:iv:eq:weyl-relacion}, le
produit \((X^aZ^b)^*X^{a'}Z^{b'}\) est un facteur de module un
multiplié par \(X^{a'-a}Z^{b'-b}\). Par conséquent,
\[
 \tau_1\bigl((X^aZ^b)^*X^{a'}Z^{b'}\bigr)
 =\delta_{a,a'}\delta_{b,b'}.
\]
Il s'agit de \(81\) matrices linéairement indépendantes dans un espace de dimension \(81\). Elles engendrent toute l'algèbre matricielle, dont le commutant est formé des scalaires ; il en résulte l'irréductibilité.
\end{proof}

\begin{theorem}[Unicité à caractère central primitif fixé]
\label{x_reciprocidad:iv:thm:weyl-unicidad}
Toute représentation unitaire irréductible du groupe présenté par
\(X^9=Z^9=c^9=I\), \(c\) central et \(ZX=cXZ\), dans laquelle
\(c\) agit comme \(\omega I\), est unitairement équivalente à
\eqref{x_reciprocidad:iv:eq:weyl-operadores}.
\end{theorem}
\begin{proof}
Soit \(v\neq0\) un vecteur propre de \(Z\), de valeur propre \(\lambda\). Comme \(Z^9=I\), \(\lambda\) est une puissance de \(\omega\). Les vecteurs \(X^kv\), \(0\leq k<9\), ont pour valeurs propres \(\omega^k\lambda\), distinctes car le caractère est primitif. Ils sont orthogonaux, et l'espace qu'ils engendrent est invariant par \(X\), \(Z\) et le centre. L'irréductibilité impose qu'il soit l'espace entier. En réordonnant la base pour commencer à la valeur propre \(1\), \(Z\) est diagonal d'entrées \(1,\omega,\ldots,\omega^8\), et \(X\) est la translation cyclique. La normalisation de la base produit l'équivalence unitaire affirmée.
\end{proof}

La primitivité du caractère et l’irréductibilité font partie de
l’énoncé. À la profondeur \(d\), les opérateurs agissant sur des facteurs
distincts commutent, et les produits tensoriels des bases précédentes
engendrent \(\mathcal A_d\). La structure de phase est ainsi
représentée à chaque niveau sans identifier le retour de phase au
retour de l’état généalogique complet.

\subsection{Deux relèvements et deux normalisations différentes}

La factorisation \eqref{x_reciprocidad:iv:eq:factorizacion-hilbert} permet de conserver
toute la nouvelle fibre ou de fixer l’une de ses étiquettes. Pour
\(p_j=|\delta_j\rangle\langle\delta_j|\), nous définissons
\begin{equation}
 \iota_d(a)=a\otimes I_9,
 \qquad \jmath_{d,j}(a)=a\otimes p_j.
 \label{x_reciprocidad:iv:eq:elevaciones}
\end{equation}
\begin{proposition}[Conservation de l’unité et de la trace]
\label{x_reciprocidad:iv:prop:unidad-traza}
L’inclusion \(\iota_d\) est unitale et conserve la trace normalisée.
L’inclusion \(\jmath_{d,j}\) est une inclusion dans un coin et
satisfait
\[
 \tau_{d+1}(\iota_d(a))=\tau_d(a),\quad \iota_d(I)=I,
 \qquad
 \tau_{d+1}(\jmath_{d,j}(a))=\frac{\tau_d(a)}9,
 \quad \jmath_{d,j}(I)=I\otimes p_j.
\]
\end{proposition}
\begin{proof}
La trace d’un produit tensoriel est le produit des traces.
Comme \(\operatorname{Tr}(I_9)=9\) et
\(\operatorname{Tr}(p_j)=1\), la division par \(9^{d+1}\)
produit les deux formules. Les images de l’identité se lisent
directement dans \eqref{x_reciprocidad:iv:eq:elevaciones}.
\end{proof}

L’inclusion unitale conserve l’ensemble des \(9\) phases et sa
normalisation globale. Le coin conserve une continuation marquée et
son projecteur de support. Les deux opérations sont légitimes, mais
transmettent des données distinctes à la profondeur suivante.

\subsection{Adhérence inductive et réalisation traciale}

\begin{theorem}[Algèbre uniformément hyperfinie nonadique]
\label{x_reciprocidad:iv:thm:uhf}
L'adhérence en norme de la réunion au moyen des inclusions \(\iota_d\) est
\[
 \mathcal A_{9^\infty}
 =\varinjlim(\mathcal A_d,\iota_d)
 \cong\bigotimes_{n=1}^{\infty}M_9(\CC).
\]
C'est une algèbre UHF de type \(9^\infty=3^\infty\), munie d'une unique trace normalisée \(\tau\), dont la restriction à \(\mathcal A_d\) est \(\tau_d\).
\end{theorem}
\begin{proof}
Les inclusions sont injectives et unitales. L’adhérence de l’union
croissante d’algèbres matricielles est précisément la construction UHF.
La compatibilité de la proposition \ref{x_reciprocidad:iv:prop:unidad-traza}
définit \(\tau\) sur l’union ; la positivité et la norme un permettent
de l’étendre à l’adhérence. Toute trace normalisée de la limite se restreint
à l’unique trace normalisée de chaque algèbre matricielle. Comme l’union
est dense, ces restrictions déterminent la trace de manière unique.
\end{proof}

Soit \(\mathcal H_\tau\) la complétion de
\(\mathcal A_{9^\infty}\) pour
\(\langle a,b\rangle_\tau=\tau(a^*b)\), et soit \(\pi_\tau\)
sa représentation par multiplication à gauche. Le vecteur défini
par l’identité est noté \(\Omega_\tau\).

\begin{theorem}[Facteur tracial du raffinement complet]
\label{x_reciprocidad:iv:thm:factor-tracial}
L'adhérence faible \(\mathcal R_9=\pi_\tau(\mathcal A_{9^\infty})''\) est un facteur hyperfini de type \(II_1\). Les inclusions nonadiques y conservent l'unité et la trace normalisée.
\end{theorem}
\begin{proof}
La trace produit se prolonge en une trace normale fidèle et finie
dans la représentation traciale. Pour voir que le centre est trivial,
soit \(E_d\) l’espérance traciale sur \(\pi_\tau(\mathcal A_d)\).
Sur un tenseur de longueur \(n>d\), elle prend la trace des \(n-d\)
derniers facteurs. Ces applications se prolongent par continuité
aux projections orthogonales correspondantes dans \(L^2\).
Si \(z\) est central, \(E_d(z)\) commute avec
\(\pi_\tau(\mathcal A_d)\), par la bimodularité de \(E_d\).
Il est donc scalaire et égal à \(\tau(z)I\). La densité de
l’union matricielle implique \(E_d(z)\to z\) dans \(L^2\), d’où
\(z=\tau(z)I\). L’algèbre est donc un facteur fini.
Elle contient unitalement des matrices de tailles \(9^d\) non bornées,
de sorte qu’elle n’est pas un facteur matriciel fini. Elle est de type \(II_1\).
Elle est hyperfinie car cette même union de matrices est faiblement
dense. La conservation de l’unité et de la trace provient de chaque inclusion
finie et se maintient lors du passage aux fermetures.
\end{proof}

\subsection{Rangs entiers et dimension traciale continue}

Dans \(\mathcal A_d\), une projection \(P\) a pour dimension normalisée \(\tau_d(P)=\operatorname{rank}(P)/9^d\). Les rangs restent entiers à toute profondeur finie. Leurs normalisations compatibles permettent le passage à la limite suivant.

\begin{proposition}[Réalisation de chaque dimension traciale]
\label{x_reciprocidad:iv:prop:dimension-tracial}
Pour chaque \(t\in[0,1]\), il existe une projection
\(P_t\in\mathcal R_9\) avec \(\tau(P_t)=t\).
\end{proposition}
\begin{proof}
Soit \(m_d=\lfloor9^dt\rfloor\). On a \(0\leq m_{d+1}-9m_d\leq8\). Une projection diagonale de rang \(m_d\) étant choisie, son inclusion a pour rang \(9m_d\) ; on ajoute \(m_{d+1}-9m_d\) positions diagonales dans son complément. On obtient ainsi une suite croissante de projections dans \(\mathcal R_9\). Sa limite forte est une projection \(P_t\), et la normalité de la trace donne
\[
 \tau(P_t)=\lim_{d\to\infty}\frac{\lfloor9^dt\rfloor}{9^d}=t.
\]
\end{proof}

Ici, \(t\) désigne la coordonnée de la dimension représentée
dans l’adhérence, et non une entrée du générateur APP. L’assertion
décrit l’étendue de la réalisation traciale résultante. Pour
la comparer à un observable déterminé, il reste nécessaire
d’identifier la projection et l’information généalogique qu’elle conserve.

Par contraste, les isométries \(V_{d,j}\xi=\xi\otimes\delta_j\) réalisent \(\jmath_{d,j}(a)=V_{d,j}aV_{d,j}^*\). À la limite d'un chemin marqué, les espaces finis sont denses. Les coins approchent tous les opérateurs de rang fini ; leur adhérence en norme est \(\mathcal K(\mathcal H_\infty)\), et leur bicommutant est \(B(\mathcal H_\infty)\), de type \(I_\infty\). Cette différence conserve le contenu opératoriel des deux choix de \eqref{x_reciprocidad:iv:eq:elevaciones} : la sélection d'une continuation et le prolongement unital de toute la fibre ont des limites différentes.

La réalisation hilbertienne, ses opérateurs de phase et la clôture
traciale sont construits dans cette section, avec leurs choix
d’inclusion explicites. Ils constituent le
support opératoire qui sera utilisé ensuite pour l’échange de
coordonnées et pour le Hamiltonien compact.
